\section{Related Work}
\label{sec:related_work}

Our proposed L-ARC framework intersects three major research domains: parameter-efficient adapter merging, dynamic model serving, and control-theoretic continuous-depth neural networks.

\subsection{Parameter-Efficient Adapters and Model Merging}
Parameter-Efficient Fine-Tuning (PEFT) methods, such as Low-Rank Adaptation (LoRA) \cite{hu2021lora}, residual adapters \cite{houlsby2019parameter, rebuffi2017learning}, and prefix-tuning \cite{li2021prefix}, have revolutionized multi-task serving. Rather than hosting multiple monolithic model instances, modern systems deploy a single pre-trained base model alongside small task-specific adapters. 
To combine the capabilities of multiple models, weight-space merging techniques like Model Soups \cite{wortsman2022model}, TIES-Merging \cite{yadav2023ties}, DARE \cite{jin2023collaborative}, and task arithmetic \cite{ilharco2022editing, ortiz2023task} combine parameter weights offline. However, static weight merging often suffers from parameter interference and is unable to adapt dynamically to incoming inputs. L-ARC operates in the activation space, dynamically blending experts layer-by-layer without static parameter degradation.

\subsection{Dynamic Model Serving and Routing}
Dynamic routing and ensembling serve as powerful paradigms to handle heterogeneous workloads in real time. Mixture-of-Experts (MoE) architectures, such as Sparsely-Gated MoE \cite{shazeer2017outrageously}, Switch Transformers \cite{fedus2022switch}, and GShard \cite{lepikhin2020gshard}, dynamically select a subset of sub-networks per token using parametric routers \cite{zhou2022mixture, rosenbaum2017routing}. While effective, training MoE routers from scratch is computationally expensive and prone to routing collapse \cite{zoph2022st}.

For pre-existing adapters, training-free dynamic serving frameworks have recently emerged. Stateless approaches like SPS-ZCA \cite{spszca2025}, PAC-ZCA \cite{paczca2025}, and SABLE \cite{sable2025} route inputs on-the-fly based on early-stage representations. While SABLE represents a strong empirical baseline for activation blending, these stateless methods re-compute ensembling weights independently at each layer, ignoring representational dependencies across depth. This lack of memory causes high-frequency ensembling weight fluctuations (routing jitter), degrading feature consistency and resulting in fragile serving dynamics.

\subsection{Continuous-Depth Neural Networks and Control Theory}
Modeling deep neural networks as discretized continuous-depth dynamical systems has gained significant traction since the introduction of Neural Ordinary Differential Equations (Neural ODEs) \cite{chen2018neural}. This perspective interprets ResNets as Euler discretizations of continuous ODE trajectories \cite{he2016deep, weinan2017proposal, lu2018beyond}. 
Control theory and dynamical systems analysis have subsequently been employed to analyze and stabilize deep network training and representation flow \cite{haber2017stable, dupont2019augmented, massaroli2020differentiable}. For instance, Lyapunov stability and contraction mapping theory have been used to analyze gradient propagation and prove convergence guarantees in deep architectures \cite{banburski2019lyapunov, jentzen2020mathematical, jordan2020lyapunov}.

ChemMerge \cite{chemmerge2025} was the first to merge physical reaction kinetics with continuous-depth ensembling, modeling task experts as chemical species reacting along layers. Its continuous ODE formulation acts as an elegant spatial low-pass filter, resolving SABLE's high-frequency routing jitter. However, ChemMerge's active representation coupling is open-loop and heuristic. Setting a constant, non-zero feedback step size $\eta$ causes cascading representational drift under heterogeneous workloads, rendering active coupling unusable. L-ARC resolves this fundamental limitation by introducing a closed-loop Lyapunov feedback controller, mathematically guaranteeing feature stability and permitting active representation coupling under any streaming conditions.
